% =============================================================================
% SECTION 2: THE ARCHITECTURAL CLASS F
% =============================================================================
\section{The Architectural Class $\cF$}\label{sec:class-f}

The Universality Theorem applies to models belonging to an architectural class $\cF$ defined by three conditions. These conditions are not merely technical regularity assumptions; each operationalizes a substantive property that distinguishes foundation MLIPs from earlier generations of machine-learned potentials.

\subsection{Condition (F1): Pretraining on Chemically Diverse DFT Corpora at Scale}

Condition (F1) requires that the model $M$ be pretrained on a DFT corpus satisfying coverage, diversity, and scale criteria. Formally, let $\cD_{\text{train}}$ denote the training dataset and $\cX$ the space of atomic configurations. Then:
\begin{enumerate}[label=(F1\alph*)]
    \item \textbf{Coverage.} The empirical distribution of $\cD_{\text{train}}$ has support intersecting every chemically relevant basin of the PES for the element combinations represented, including at minimum equilibrium structures, strained configurations, and surface terminations.
    \item \textbf{Diversity.} The training set spans at least 80\% of the stable element combinations in the periodic table (as catalogued in the Materials Project\cite{jain2013materials} or Alexandria database\cite{schmidt2024alexandria}).
    \item \textbf{Scale.} The number of distinct DFT calculations exceeds $10^6$, with total atom-instances exceeding $10^8$.
\end{enumerate}

Published corpora satisfying (F1) include the Materials Project trajectory data (MPtrj)\cite{jain2013materials}, the Alexandria/sAlex dataset\cite{schmidt2024alexandria}, OMat24\cite{barroso2024omat24}, and the MAD dataset\cite{mazitov2025petmad}. The OMat24 dataset, for instance, comprises ``2 to 3-fold more distinct atomic environments across the entire periodic table compared to previous databases''\cite{yang2024mattersim}. Condition (F1) is the empirical anchor for the uniform training-distribution coverage constant $\kappa_1$ that appears in the Cross-Model Vandermonde Lemma.

\subsection{Condition (F2): Representational Expressivity Above the Incompleteness Floor}

Condition (F2) requires that the model architecture escape the Pozdnyakov--Ceriotti incompleteness floor\cite{pozdnyakov2020incompleteness}. Pozdnyakov \emph{et al.}\ proved that three- and four-body invariant descriptors cannot uniquely distinguish all non-isomorphic atomic environments: there exist distinct configurations with identical three-body and four-body correlation functions. This incompleteness implies that any MLIP whose representation is limited to such descriptors must suffer systematic errors on configurations whose true physics requires higher-body-order discrimination.

The escape from this floor is operationalized via the geometric Weisfeiler--Leman (GWL) hierarchy developed by Joshi \emph{et al.}\cite{joshi2023expressive} and the completeness framework of Nigam \emph{et al.}\cite{nigam2023completeness}. Specifically, (F2) requires that $M$ employ an architecture at least as expressive as an equivariant message-passing network of depth $\geq 2$ with body-order $\geq 4$ features. In the GWL hierarchy, this means:
\begin{equation}\label{eq:f2-gwl}
    \text{GWL}(M) \geq \text{GWL}_{\text{equivariant}}^{(2,4)},
\end{equation}
where the superscript $(2,4)$ denotes depth 2 and body order 4. Architectures satisfying (F2) include MACE\cite{batatia2022mace}, NequIP\cite{batzner2022nequip}, Allegro\cite{musaelian2023allegro}, eqV2\cite{liao2024equiformerv2}, GRACE\cite{bochkarev2024grace}, eSEN\cite{om枯萎}, PET\cite{batatia2024pet}, SevenNet\cite{park2024sevennet}, Orb-v3\cite{qian2023orb}, MatterSim\cite{yang2024mattersim}, and DPA3\cite{zhang2024dpa3}. The margin by which $M$ exceeds the incompleteness floor contributes to the constant $\kappa_2$ in the Vandermonde lemma.

\subsection{Condition (F3): Smoothness Outside a Measure-Zero Singular Set}

Condition (F3) requires that the model $M: \cX \to \R$ be a $C^2$ function of atomic positions outside a measure-zero set $\cS_M \subset \cX$ of singular configurations. The singular set typically contains:
\begin{itemize}
    \item Coincident-atom configurations (nuclear coalescence points);
    \item Symmetry-enhanced configurations where degenerate eigenvalues of the descriptor matrix cause non-analyticity;
    \item Phase-transition points where the electronic structure changes discontinuously.
\end{itemize}

The smoothness requirement is anchored to the smooth basis construction of Bigi \emph{et al.}\cite{bigi2022smooth}, who proved that equivariant message-passing networks with smooth radial bases and smooth activation functions yield $C^\infty$ PES predictions away from coincident-atom singularities. The standard equivariant-message-passing architectures (MACE, NequIP, Allegro) satisfy this condition with a Lipschitz constant $L_M$ on the Jacobian $\Jac M(x)$ that is bounded uniformly on compact sublevel sets of the loss landscape. The class-uniform smoothness constant $\kappa_3$ in the Vandermonde lemma is the supremum of $L_M$ over $M \in \cF$.

\subsection{Worked Examples and Boundary Cases}

Table~\ref{tab:class-f-examples} classifies representative models according to (F1)--(F3).

\begin{table}[htbp]
\centering
\caption{Classification of representative MLIP architectures with respect to conditions (F1)--(F3).}
\label{tab:class-f-examples}
\begin{tabular}{@{}lccccl@{}}
\toprule
\textbf{Model} & \textbf{F1} & \textbf{F2} & \textbf{F3} & $\in \cF$ & \textbf{Notes} \\
\midrule
MACE-MP-0/MPA-0 & Yes & Yes & Yes & Yes & E(3)-equivariant, MP-traj + sAlex \\
CHGNet & Yes & Marginal & Yes & Yes* & GWL depth 2, marginal body order \\
Orb-v3 & Yes & Yes & Yes & Yes & Equivariant transformer \\
MatterSim & Yes & Yes & Yes & Yes & 0--5000~K, 0--1000~GPa training \\
eqV2/UMA & Yes & Yes & Yes & Yes & Equiformer architecture \\
PET-MAD & Yes & Yes & Yes & Yes & Broadened training distribution \\
SevenNet & Yes & Yes & Yes & Yes & GNN with scalar embedding \\
GRACE & Yes & Yes & Yes & Yes & Graph ACE, semilocal interactions \\
eSEN & Yes & Yes & Yes & Yes & Equivariant spherical harmonics \\
DPA3 & Yes & Yes & Yes & Yes & Attention-based, DPA-1/2 successor \\
M3GNet & Yes & Marginal & Yes & Yes* & Earlier generation, marginal F2 \\
\midrule
Classical EAM & N/A & No & Yes & No & Fails F2 (incomplete descriptor) \\
GAP-original & Partial & No & Yes & No & Fails F2 at scale (SOAP only) \\
SchNet & Yes & No & Yes & No & Fails F2 (invariant, depth 1) \\
\bottomrule
\end{tabular}
\end{table}

Models marked Yes* (CHGNet, M3GNet) are boundary cases: they marginally satisfy (F2) at the minimum GWL depth and body order, and their inclusion in $\cF$ depends on whether the specific instantiation achieves the expressivity threshold in practice. The operational test is whether the model's descriptor can distinguish the counterexample configurations constructed by Pozdnyakov \emph{et al.}\cite{pozdnyakov2020incompleteness}; CHGNet-v2 and later variants pass this test, while the original M3GNet does not.

\subsection{Closure Properties}

The class $\cF$ is closed under two operations relevant to the theorem:
\begin{enumerate}
    \item \textbf{Fine-tuning.} If $M \in \cF$ and $M'$ is obtained from $M$ by gradient-based fine-tuning on a downstream dataset, then $M' \in \cF$ provided the fine-tuning does not alter the architecture (hence F2 and F3 are preserved) and the combined training corpus still satisfies (F1). This closure property underlies clause (vi) on generational stability.
    \item \textbf{Ensemble averaging.} If $M_1, M_2 \in \cF$ and $M_{\text{ens}} = \sum_i w_i M_i$ with $\sum_i w_i = 1$, then $M_{\text{ens}} \in \cF$ provided each component satisfies (F1)--(F3) individually. The smoothness constant $\kappa_3$ for the ensemble is bounded by $\max_i \kappa_3(M_i)$.
\end{enumerate}
